\section{Introduction}
\label{introduction}

\begin{figure}[htbp]
\centering
\includegraphics[width=0.6\linewidth, height=0.4\linewidth]{page1-image-963-s15.png}  
\caption{\textbf{Top Row}: A corresponding dynamic-static LiDAR pair $(D,S)$. Left - Dynamic scan, $D$ which has dynamic objects (in \textcolor{blue}{blue}) and dynamic occlusions (indicated by arrows). Right - corresponding static scan, $S$ in the same environment. \textbf{Bottom Row}: Our model in action. Given the dynamic scan $D$ as input, our model generates a corresponding static scan $(\hat{S})$ (middle), that approximates $S$. We subtract $D$ from $\hat{S}$  to extract moving and movable objects (pink boxes in rightmost image). }
\label{fig:movable-first}
\end{figure}

Handling non-stationary objects is an important problem with major applications for LiDAR (Light Detection and Ranging) based Autonomous Vehicles (AVs). Current solutions extract existing static structures by segmenting and removing moving objects  (\citet{chen2021moving}) with the help of segmentation and object information. They use this information to train a model to identify and remove moving objects in real time. Collecting segmentation labels is laborious, expensive and time-consuming, especially for several popular and important settings like industrial environments, fulfilment centers, urban settings, etc. In addition, common everyday scenarios also involve \textit{movable} objects. They have the capacity to move but may be stationary at a given instant, and can escape detection in segmentation label free settings due to uncertain motion. When moving, these objects add to the moving parts of a scene, thereby leading to erroneous pose estimates.  Thus, they are a potential hazard to AVs.

We address the above issues by proposing \dslr{}, a GAN based generative network that enables precise segmentation of dynamic - movable as well as moving objects. This is achieved in the absence of any form of segmentation information, by generating corresponding static reconstructions for a given dynamic LiDAR scan (Figure \ref{fig:movable-first}). \dslr{} is trained with corresponding static and dynamic LiDAR scan pairs which are easier to collect and are time and cost effective (\citet{dslr-git}) as compared to labelled information collection. In \dslr{}, we introduce an intuitively designed discriminator that enforces an implicit attention to dynamic objects. This is coupled with an intuitively designed contrastive strategy that reinforces the dissimilarity between dynamic and static scans. 

We list the contributions of our work here:
\begin{itemize}

    \item We present \dslr{}, a GAN-based adversarial model for \textit{generating static background} of occluded parts of a dynamic LiDAR scan (dynamic to static translation) and henceforth \textit{segmenting out moving and movable  parts of scenes}. This is achieved in a segmentation-label devoid setting. \dslr{} uses a novel couple discriminator to discriminate carefully chosen pairs of static and dynamic input. This design choice helps our model to explicitly focus its attention on moving as well as movable objects, without external supervision.
 
    \item To reinforce the attention of our couple discriminator on the pattern of moving and movable objects, we design a novel contrastive strategy. We utilize static-dynamic LiDAR correspondence to design hard negative triplets of scans for contrastive learning. 
    
    \item We quantitatively demonstrate the superior quality of static reconstructions on 3 varied datasets. We qualitatively demonstrate the usefulness of \dslr{} in segmenting out movable as well as moving objects without using segmentation labels. 
       \item To work well with datasets lacking static-dynamic correspondence, we propose \dslruda{}. We use Unsupervised Domain Adaptation to adapt models trained on correspondence datasets to datasets lacking correspondence. Despite working in segmentation devoid settings, \dslruda{} performs comparable to or better than segmentation based baselines for navigation in long sequences with \textit{highly dynamic environments}.
\end{itemize}